\documentclass[a4paper,12pt,titlepage]{article}

%Εισαγωγή γλωσσικής υποστήριξης
%ελληνικό hyphenation

\usepackage{fontspec}
\usepackage{xunicode}
\usepackage{xltxtra}
\usepackage{xgreek}
\usepackage[colorlinks]{hyperref}
\usepackage{enumerate}
\usepackage{amsmath}
%η γραμματοσειρά
%\setmainfont[Mapping=tex-text]{Times New Roman} %απλοποιημένο σε σχέση με το άρθρο
\setmainfont[Mapping=tex-text]{Linux Libertine} 
%page orientation
\usepackage{a4wide}
\voffset = -0.5in
\textheight = 664pt

%Άλλα χρήσιμα πακέτα
\usepackage{graphicx} 	%εισαγωγή εικόνων jpg/png κλπ
\usepackage{listings}
\lstset{commentstyle=\textit,
captionpos=b,
breakatwhitespace=true,
showstringspaces=false,
breaklines=true,
keywordstyle=\color{black}\bfseries,
float=htb,
frame=single}

%απενεργοποίηση του indent στις νέες παραγράφους
\parindent=0in

%macro που δίνει το μέγιστο επιτρεπτό μέγεθος σε μια εικόνα χωρίς να παραβιάζει τα όρια του LaTeX
\makeatletter
\def\maxwidth{
\ifdim\Gin@nat@width>\linewidth
\linewidth
\else
\Gin@nat@width
\fi
}

\makeatother

\title{1η Άσκηση\\kn}
\author{AxiP}
\date{\today}

\begin{document}

\pagestyle{headings}    %αρίθμηση στο πάνω μέρος της σελίδας

%\maketitle
\begin{titlepage}
\begin{center}
\includegraphics[width=50mm]{pyrforos.pdf}\\[0.5cm]
\textbf{\LARGE ΕΘΝΙΚΟ ΜΕΤΣΟΒΙΟ ΠΟΛΥΤΕΧΝΕΙΟ}\\
\textrm{\Large Σχολή Εφαρμοσμένων Μαθηματικών και Φυσικών Επιστημών}\\[2.0cm]
\Huge{Τεχνολογία των laser\\}
\Large{\textit{6o εξάμηνο, ΣΕΜΦΕ}}\\[2.0cm]
\Large{\textit{\textbf{Κατασκευή τροφοδοτικού διοδικού laser}}}\\[5.0cm]
\normalsize

\begin{minipage}{0.49\textwidth}
\begin{flushleft}
\textbf{Αχιλλέας Πιπινέλλης}, 09103163
\end{flushleft}
\end{minipage}
\begin{minipage}{0.49\textwidth}
\begin{flushright}
\textbf{\\
\textit{Η/Μ Παράδοσης:} 2 Ιουνίου 2012}
\end{flushright}
\end{minipage}

%\maketitle

\vfill
%bottom of the page
{Αθήνα, 2011}

\end{center}
\end{titlepage}

\section{Ερώτημα 1}

\begin{center}
% GNUPLOT: LaTeX picture
\setlength{\unitlength}{0.240900pt}
\ifx\plotpoint\undefined\newsavebox{\plotpoint}\fi
\sbox{\plotpoint}{\rule[-0.200pt]{0.400pt}{0.400pt}}%
\begin{picture}(1500,900)(0,0)
\sbox{\plotpoint}{\rule[-0.200pt]{0.400pt}{0.400pt}}%
\put(191.0,131.0){\rule[-0.200pt]{4.818pt}{0.400pt}}
\put(171,131){\makebox(0,0)[r]{$-0.1$}}
\put(1419.0,131.0){\rule[-0.200pt]{4.818pt}{0.400pt}}
\put(191.0,212.0){\rule[-0.200pt]{4.818pt}{0.400pt}}
\put(171,212){\makebox(0,0)[r]{$-0.05$}}
\put(1419.0,212.0){\rule[-0.200pt]{4.818pt}{0.400pt}}
\put(191.0,292.0){\rule[-0.200pt]{4.818pt}{0.400pt}}
\put(171,292){\makebox(0,0)[r]{$0$}}
\put(1419.0,292.0){\rule[-0.200pt]{4.818pt}{0.400pt}}
\put(191.0,373.0){\rule[-0.200pt]{4.818pt}{0.400pt}}
\put(171,373){\makebox(0,0)[r]{$0.05$}}
\put(1419.0,373.0){\rule[-0.200pt]{4.818pt}{0.400pt}}
\put(191.0,454.0){\rule[-0.200pt]{4.818pt}{0.400pt}}
\put(171,454){\makebox(0,0)[r]{$0.1$}}
\put(1419.0,454.0){\rule[-0.200pt]{4.818pt}{0.400pt}}
\put(191.0,534.0){\rule[-0.200pt]{4.818pt}{0.400pt}}
\put(171,534){\makebox(0,0)[r]{$0.15$}}
\put(1419.0,534.0){\rule[-0.200pt]{4.818pt}{0.400pt}}
\put(191.0,615.0){\rule[-0.200pt]{4.818pt}{0.400pt}}
\put(171,615){\makebox(0,0)[r]{$0.2$}}
\put(1419.0,615.0){\rule[-0.200pt]{4.818pt}{0.400pt}}
\put(191.0,695.0){\rule[-0.200pt]{4.818pt}{0.400pt}}
\put(171,695){\makebox(0,0)[r]{$0.25$}}
\put(1419.0,695.0){\rule[-0.200pt]{4.818pt}{0.400pt}}
\put(191.0,776.0){\rule[-0.200pt]{4.818pt}{0.400pt}}
\put(171,776){\makebox(0,0)[r]{$0.3$}}
\put(1419.0,776.0){\rule[-0.200pt]{4.818pt}{0.400pt}}
\put(191.0,131.0){\rule[-0.200pt]{0.400pt}{4.818pt}}
\put(191,90){\makebox(0,0){$0$}}
\put(191.0,756.0){\rule[-0.200pt]{0.400pt}{4.818pt}}
\put(399.0,131.0){\rule[-0.200pt]{0.400pt}{4.818pt}}
\put(399,90){\makebox(0,0){$200$}}
\put(399.0,756.0){\rule[-0.200pt]{0.400pt}{4.818pt}}
\put(607.0,131.0){\rule[-0.200pt]{0.400pt}{4.818pt}}
\put(607,90){\makebox(0,0){$400$}}
\put(607.0,756.0){\rule[-0.200pt]{0.400pt}{4.818pt}}
\put(815.0,131.0){\rule[-0.200pt]{0.400pt}{4.818pt}}
\put(815,90){\makebox(0,0){$600$}}
\put(815.0,756.0){\rule[-0.200pt]{0.400pt}{4.818pt}}
\put(1023.0,131.0){\rule[-0.200pt]{0.400pt}{4.818pt}}
\put(1023,90){\makebox(0,0){$800$}}
\put(1023.0,756.0){\rule[-0.200pt]{0.400pt}{4.818pt}}
\put(1231.0,131.0){\rule[-0.200pt]{0.400pt}{4.818pt}}
\put(1231,90){\makebox(0,0){$1000$}}
\put(1231.0,756.0){\rule[-0.200pt]{0.400pt}{4.818pt}}
\put(1439.0,131.0){\rule[-0.200pt]{0.400pt}{4.818pt}}
\put(1439,90){\makebox(0,0){$1200$}}
\put(1439.0,756.0){\rule[-0.200pt]{0.400pt}{4.818pt}}
\put(191.0,131.0){\rule[-0.200pt]{0.400pt}{155.380pt}}
\put(191.0,131.0){\rule[-0.200pt]{300.643pt}{0.400pt}}
\put(1439.0,131.0){\rule[-0.200pt]{0.400pt}{155.380pt}}
\put(191.0,776.0){\rule[-0.200pt]{300.643pt}{0.400pt}}
\put(30,453){\makebox(0,0){$I/I_0$}}
\put(815,29){\makebox(0,0){$U(Volts)$}}
\put(815,838){\makebox(0,0){$I/I_0 (U)$}}
\put(295,308){\usebox{\plotpoint}}
\put(295,308.17){\rule{2.500pt}{0.400pt}}
\multiput(295.00,307.17)(6.811,2.000){2}{\rule{1.250pt}{0.400pt}}
\put(307,308.67){\rule{2.650pt}{0.400pt}}
\multiput(307.00,309.17)(5.500,-1.000){2}{\rule{1.325pt}{0.400pt}}
\multiput(318.00,307.93)(1.033,-0.482){9}{\rule{0.900pt}{0.116pt}}
\multiput(318.00,308.17)(10.132,-6.000){2}{\rule{0.450pt}{0.400pt}}
\multiput(330.00,301.93)(0.692,-0.488){13}{\rule{0.650pt}{0.117pt}}
\multiput(330.00,302.17)(9.651,-8.000){2}{\rule{0.325pt}{0.400pt}}
\multiput(341.00,293.94)(1.651,-0.468){5}{\rule{1.300pt}{0.113pt}}
\multiput(341.00,294.17)(9.302,-4.000){2}{\rule{0.650pt}{0.400pt}}
\put(353,290.67){\rule{2.650pt}{0.400pt}}
\multiput(353.00,290.17)(5.500,1.000){2}{\rule{1.325pt}{0.400pt}}
\put(376,290.17){\rule{2.300pt}{0.400pt}}
\multiput(376.00,291.17)(6.226,-2.000){2}{\rule{1.150pt}{0.400pt}}
\put(387,290.17){\rule{2.500pt}{0.400pt}}
\multiput(387.00,289.17)(6.811,2.000){2}{\rule{1.250pt}{0.400pt}}
\multiput(399.00,292.59)(0.669,0.489){15}{\rule{0.633pt}{0.118pt}}
\multiput(399.00,291.17)(10.685,9.000){2}{\rule{0.317pt}{0.400pt}}
\multiput(411.00,301.59)(0.798,0.485){11}{\rule{0.729pt}{0.117pt}}
\multiput(411.00,300.17)(9.488,7.000){2}{\rule{0.364pt}{0.400pt}}
\multiput(422.00,306.95)(2.472,-0.447){3}{\rule{1.700pt}{0.108pt}}
\multiput(422.00,307.17)(8.472,-3.000){2}{\rule{0.850pt}{0.400pt}}
\multiput(434.00,303.93)(0.611,-0.489){15}{\rule{0.589pt}{0.118pt}}
\multiput(434.00,304.17)(9.778,-9.000){2}{\rule{0.294pt}{0.400pt}}
\multiput(445.00,294.93)(1.033,-0.482){9}{\rule{0.900pt}{0.116pt}}
\multiput(445.00,295.17)(10.132,-6.000){2}{\rule{0.450pt}{0.400pt}}
\put(457,289.67){\rule{2.650pt}{0.400pt}}
\multiput(457.00,289.17)(5.500,1.000){2}{\rule{1.325pt}{0.400pt}}
\put(468,291.17){\rule{2.500pt}{0.400pt}}
\multiput(468.00,290.17)(6.811,2.000){2}{\rule{1.250pt}{0.400pt}}
\put(364.0,292.0){\rule[-0.200pt]{2.891pt}{0.400pt}}
\put(491,291.67){\rule{2.891pt}{0.400pt}}
\multiput(491.00,292.17)(6.000,-1.000){2}{\rule{1.445pt}{0.400pt}}
\put(503,290.67){\rule{2.891pt}{0.400pt}}
\multiput(503.00,291.17)(6.000,-1.000){2}{\rule{1.445pt}{0.400pt}}
\put(480.0,293.0){\rule[-0.200pt]{2.650pt}{0.400pt}}
\multiput(526.00,291.59)(1.033,0.482){9}{\rule{0.900pt}{0.116pt}}
\multiput(526.00,290.17)(10.132,6.000){2}{\rule{0.450pt}{0.400pt}}
\multiput(538.00,297.59)(0.692,0.488){13}{\rule{0.650pt}{0.117pt}}
\multiput(538.00,296.17)(9.651,8.000){2}{\rule{0.325pt}{0.400pt}}
\multiput(549.00,305.59)(1.267,0.477){7}{\rule{1.060pt}{0.115pt}}
\multiput(549.00,304.17)(9.800,5.000){2}{\rule{0.530pt}{0.400pt}}
\put(515.0,291.0){\rule[-0.200pt]{2.650pt}{0.400pt}}
\put(572,308.17){\rule{2.500pt}{0.400pt}}
\multiput(572.00,309.17)(6.811,-2.000){2}{\rule{1.250pt}{0.400pt}}
\put(584,306.17){\rule{2.300pt}{0.400pt}}
\multiput(584.00,307.17)(6.226,-2.000){2}{\rule{1.150pt}{0.400pt}}
\put(595,306.17){\rule{2.500pt}{0.400pt}}
\multiput(595.00,305.17)(6.811,2.000){2}{\rule{1.250pt}{0.400pt}}
\multiput(607.00,308.59)(0.874,0.485){11}{\rule{0.786pt}{0.117pt}}
\multiput(607.00,307.17)(10.369,7.000){2}{\rule{0.393pt}{0.400pt}}
\multiput(619.00,315.59)(0.692,0.488){13}{\rule{0.650pt}{0.117pt}}
\multiput(619.00,314.17)(9.651,8.000){2}{\rule{0.325pt}{0.400pt}}
\multiput(630.00,323.61)(2.472,0.447){3}{\rule{1.700pt}{0.108pt}}
\multiput(630.00,322.17)(8.472,3.000){2}{\rule{0.850pt}{0.400pt}}
\put(561.0,310.0){\rule[-0.200pt]{2.650pt}{0.400pt}}
\multiput(653.00,324.95)(2.472,-0.447){3}{\rule{1.700pt}{0.108pt}}
\multiput(653.00,325.17)(8.472,-3.000){2}{\rule{0.850pt}{0.400pt}}
\put(665,321.67){\rule{2.650pt}{0.400pt}}
\multiput(665.00,322.17)(5.500,-1.000){2}{\rule{1.325pt}{0.400pt}}
\multiput(676.00,322.60)(1.651,0.468){5}{\rule{1.300pt}{0.113pt}}
\multiput(676.00,321.17)(9.302,4.000){2}{\rule{0.650pt}{0.400pt}}
\multiput(688.00,326.59)(0.692,0.488){13}{\rule{0.650pt}{0.117pt}}
\multiput(688.00,325.17)(9.651,8.000){2}{\rule{0.325pt}{0.400pt}}
\multiput(699.00,334.59)(0.874,0.485){11}{\rule{0.786pt}{0.117pt}}
\multiput(699.00,333.17)(10.369,7.000){2}{\rule{0.393pt}{0.400pt}}
\put(711,340.67){\rule{2.891pt}{0.400pt}}
\multiput(711.00,340.17)(6.000,1.000){2}{\rule{1.445pt}{0.400pt}}
\put(723,340.67){\rule{2.650pt}{0.400pt}}
\multiput(723.00,341.17)(5.500,-1.000){2}{\rule{1.325pt}{0.400pt}}
\put(734,339.67){\rule{2.891pt}{0.400pt}}
\multiput(734.00,340.17)(6.000,-1.000){2}{\rule{1.445pt}{0.400pt}}
\put(642.0,326.0){\rule[-0.200pt]{2.650pt}{0.400pt}}
\put(757,339.67){\rule{2.891pt}{0.400pt}}
\multiput(757.00,339.17)(6.000,1.000){2}{\rule{1.445pt}{0.400pt}}
\put(769,340.67){\rule{2.650pt}{0.400pt}}
\multiput(769.00,340.17)(5.500,1.000){2}{\rule{1.325pt}{0.400pt}}
\put(780,340.17){\rule{2.500pt}{0.400pt}}
\multiput(780.00,341.17)(6.811,-2.000){2}{\rule{1.250pt}{0.400pt}}
\multiput(792.00,338.95)(2.248,-0.447){3}{\rule{1.567pt}{0.108pt}}
\multiput(792.00,339.17)(7.748,-3.000){2}{\rule{0.783pt}{0.400pt}}
\multiput(803.00,337.60)(1.651,0.468){5}{\rule{1.300pt}{0.113pt}}
\multiput(803.00,336.17)(9.302,4.000){2}{\rule{0.650pt}{0.400pt}}
\multiput(815.58,341.00)(0.492,0.712){21}{\rule{0.119pt}{0.667pt}}
\multiput(814.17,341.00)(12.000,15.616){2}{\rule{0.400pt}{0.333pt}}
\multiput(827.58,358.00)(0.492,1.156){19}{\rule{0.118pt}{1.009pt}}
\multiput(826.17,358.00)(11.000,22.906){2}{\rule{0.400pt}{0.505pt}}
\multiput(838.58,383.00)(0.492,1.013){21}{\rule{0.119pt}{0.900pt}}
\multiput(837.17,383.00)(12.000,22.132){2}{\rule{0.400pt}{0.450pt}}
\multiput(850.58,407.00)(0.492,0.967){19}{\rule{0.118pt}{0.864pt}}
\multiput(849.17,407.00)(11.000,19.207){2}{\rule{0.400pt}{0.432pt}}
\multiput(861.58,428.00)(0.492,0.755){21}{\rule{0.119pt}{0.700pt}}
\multiput(860.17,428.00)(12.000,16.547){2}{\rule{0.400pt}{0.350pt}}
\multiput(873.58,446.00)(0.492,0.826){19}{\rule{0.118pt}{0.755pt}}
\multiput(872.17,446.00)(11.000,16.434){2}{\rule{0.400pt}{0.377pt}}
\multiput(884.58,464.00)(0.492,1.315){21}{\rule{0.119pt}{1.133pt}}
\multiput(883.17,464.00)(12.000,28.648){2}{\rule{0.400pt}{0.567pt}}
\multiput(896.58,495.00)(0.492,2.100){19}{\rule{0.118pt}{1.736pt}}
\multiput(895.17,495.00)(11.000,41.396){2}{\rule{0.400pt}{0.868pt}}
\multiput(907.58,540.00)(0.492,1.789){21}{\rule{0.119pt}{1.500pt}}
\multiput(906.17,540.00)(12.000,38.887){2}{\rule{0.400pt}{0.750pt}}
\multiput(919.58,582.00)(0.492,1.229){21}{\rule{0.119pt}{1.067pt}}
\multiput(918.17,582.00)(12.000,26.786){2}{\rule{0.400pt}{0.533pt}}
\multiput(931.58,611.00)(0.492,0.732){19}{\rule{0.118pt}{0.682pt}}
\multiput(930.17,611.00)(11.000,14.585){2}{\rule{0.400pt}{0.341pt}}
\multiput(942.58,627.00)(0.492,0.582){21}{\rule{0.119pt}{0.567pt}}
\multiput(941.17,627.00)(12.000,12.824){2}{\rule{0.400pt}{0.283pt}}
\multiput(954.58,641.00)(0.492,0.637){19}{\rule{0.118pt}{0.609pt}}
\multiput(953.17,641.00)(11.000,12.736){2}{\rule{0.400pt}{0.305pt}}
\multiput(965.58,655.00)(0.492,0.712){21}{\rule{0.119pt}{0.667pt}}
\multiput(964.17,655.00)(12.000,15.616){2}{\rule{0.400pt}{0.333pt}}
\multiput(977.58,672.00)(0.492,0.779){19}{\rule{0.118pt}{0.718pt}}
\multiput(976.17,672.00)(11.000,15.509){2}{\rule{0.400pt}{0.359pt}}
\multiput(988.00,689.59)(1.033,0.482){9}{\rule{0.900pt}{0.116pt}}
\multiput(988.00,688.17)(10.132,6.000){2}{\rule{0.450pt}{0.400pt}}
\multiput(1000.00,693.93)(0.798,-0.485){11}{\rule{0.729pt}{0.117pt}}
\multiput(1000.00,694.17)(9.488,-7.000){2}{\rule{0.364pt}{0.400pt}}
\multiput(1011.00,686.93)(0.669,-0.489){15}{\rule{0.633pt}{0.118pt}}
\multiput(1011.00,687.17)(10.685,-9.000){2}{\rule{0.317pt}{0.400pt}}
\multiput(1023.00,677.93)(1.267,-0.477){7}{\rule{1.060pt}{0.115pt}}
\multiput(1023.00,678.17)(9.800,-5.000){2}{\rule{0.530pt}{0.400pt}}
\multiput(1035.58,669.36)(0.492,-1.298){19}{\rule{0.118pt}{1.118pt}}
\multiput(1034.17,671.68)(11.000,-25.679){2}{\rule{0.400pt}{0.559pt}}
\multiput(1046.58,634.38)(0.492,-3.469){21}{\rule{0.119pt}{2.800pt}}
\multiput(1045.17,640.19)(12.000,-75.188){2}{\rule{0.400pt}{1.400pt}}
\multiput(1058.58,549.04)(0.492,-4.835){19}{\rule{0.118pt}{3.845pt}}
\multiput(1057.17,557.02)(11.000,-95.019){2}{\rule{0.400pt}{1.923pt}}
\multiput(1069.58,452.59)(0.492,-2.780){21}{\rule{0.119pt}{2.267pt}}
\multiput(1068.17,457.30)(12.000,-60.295){2}{\rule{0.400pt}{1.133pt}}
\multiput(1081.00,395.93)(0.692,-0.488){13}{\rule{0.650pt}{0.117pt}}
\multiput(1081.00,396.17)(9.651,-8.000){2}{\rule{0.325pt}{0.400pt}}
\multiput(1092.58,389.00)(0.492,1.013){21}{\rule{0.119pt}{0.900pt}}
\multiput(1091.17,389.00)(12.000,22.132){2}{\rule{0.400pt}{0.450pt}}
\multiput(1104.58,413.00)(0.492,1.581){19}{\rule{0.118pt}{1.336pt}}
\multiput(1103.17,413.00)(11.000,31.226){2}{\rule{0.400pt}{0.668pt}}
\multiput(1115.58,447.00)(0.492,1.659){21}{\rule{0.119pt}{1.400pt}}
\multiput(1114.17,447.00)(12.000,36.094){2}{\rule{0.400pt}{0.700pt}}
\multiput(1127.58,486.00)(0.492,1.573){21}{\rule{0.119pt}{1.333pt}}
\multiput(1126.17,486.00)(12.000,34.233){2}{\rule{0.400pt}{0.667pt}}
\multiput(1139.58,523.00)(0.492,1.628){19}{\rule{0.118pt}{1.373pt}}
\multiput(1138.17,523.00)(11.000,32.151){2}{\rule{0.400pt}{0.686pt}}
\multiput(1150.58,558.00)(0.492,1.401){21}{\rule{0.119pt}{1.200pt}}
\multiput(1149.17,558.00)(12.000,30.509){2}{\rule{0.400pt}{0.600pt}}
\multiput(1162.58,591.00)(0.492,1.298){19}{\rule{0.118pt}{1.118pt}}
\multiput(1161.17,591.00)(11.000,25.679){2}{\rule{0.400pt}{0.559pt}}
\multiput(1173.58,619.00)(0.492,0.884){21}{\rule{0.119pt}{0.800pt}}
\multiput(1172.17,619.00)(12.000,19.340){2}{\rule{0.400pt}{0.400pt}}
\multiput(1185.58,640.00)(0.492,0.684){19}{\rule{0.118pt}{0.645pt}}
\multiput(1184.17,640.00)(11.000,13.660){2}{\rule{0.400pt}{0.323pt}}
\multiput(1196.00,655.58)(0.543,0.492){19}{\rule{0.536pt}{0.118pt}}
\multiput(1196.00,654.17)(10.887,11.000){2}{\rule{0.268pt}{0.400pt}}
\multiput(1208.00,666.58)(0.547,0.491){17}{\rule{0.540pt}{0.118pt}}
\multiput(1208.00,665.17)(9.879,10.000){2}{\rule{0.270pt}{0.400pt}}
\multiput(1219.00,676.61)(2.472,0.447){3}{\rule{1.700pt}{0.108pt}}
\multiput(1219.00,675.17)(8.472,3.000){2}{\rule{0.850pt}{0.400pt}}
\multiput(1231.00,677.93)(1.033,-0.482){9}{\rule{0.900pt}{0.116pt}}
\multiput(1231.00,678.17)(10.132,-6.000){2}{\rule{0.450pt}{0.400pt}}
\multiput(1243.58,669.72)(0.492,-0.873){19}{\rule{0.118pt}{0.791pt}}
\multiput(1242.17,671.36)(11.000,-17.358){2}{\rule{0.400pt}{0.395pt}}
\multiput(1254.58,649.30)(0.492,-1.315){21}{\rule{0.119pt}{1.133pt}}
\multiput(1253.17,651.65)(12.000,-28.648){2}{\rule{0.400pt}{0.567pt}}
\multiput(1266.58,616.85)(0.492,-1.769){19}{\rule{0.118pt}{1.482pt}}
\multiput(1265.17,619.92)(11.000,-34.924){2}{\rule{0.400pt}{0.741pt}}
\multiput(1277.58,579.60)(0.492,-1.530){21}{\rule{0.119pt}{1.300pt}}
\multiput(1276.17,582.30)(12.000,-33.302){2}{\rule{0.400pt}{0.650pt}}
\multiput(1289.58,543.45)(0.492,-1.581){19}{\rule{0.118pt}{1.336pt}}
\multiput(1288.17,546.23)(11.000,-31.226){2}{\rule{0.400pt}{0.668pt}}
\multiput(1300.58,508.91)(0.492,-1.746){21}{\rule{0.119pt}{1.467pt}}
\multiput(1299.17,511.96)(12.000,-37.956){2}{\rule{0.400pt}{0.733pt}}
\multiput(1312.58,466.64)(0.492,-2.147){19}{\rule{0.118pt}{1.773pt}}
\multiput(1311.17,470.32)(11.000,-42.321){2}{\rule{0.400pt}{0.886pt}}
\multiput(1323.58,424.40)(0.492,-0.970){21}{\rule{0.119pt}{0.867pt}}
\multiput(1322.17,426.20)(12.000,-21.201){2}{\rule{0.400pt}{0.433pt}}
\multiput(1335.58,405.00)(0.492,0.841){21}{\rule{0.119pt}{0.767pt}}
\multiput(1334.17,405.00)(12.000,18.409){2}{\rule{0.400pt}{0.383pt}}
\multiput(1347.58,425.00)(0.492,2.194){19}{\rule{0.118pt}{1.809pt}}
\multiput(1346.17,425.00)(11.000,43.245){2}{\rule{0.400pt}{0.905pt}}
\multiput(1358.58,472.00)(0.492,2.176){21}{\rule{0.119pt}{1.800pt}}
\multiput(1357.17,472.00)(12.000,47.264){2}{\rule{0.400pt}{0.900pt}}
\multiput(1370.58,523.00)(0.492,1.345){19}{\rule{0.118pt}{1.155pt}}
\multiput(1369.17,523.00)(11.000,26.604){2}{\rule{0.400pt}{0.577pt}}
\multiput(1381.58,549.23)(0.492,-0.712){21}{\rule{0.119pt}{0.667pt}}
\multiput(1380.17,550.62)(12.000,-15.616){2}{\rule{0.400pt}{0.333pt}}
\multiput(1393.58,526.43)(0.492,-2.524){19}{\rule{0.118pt}{2.064pt}}
\multiput(1392.17,530.72)(11.000,-49.717){2}{\rule{0.400pt}{1.032pt}}
\multiput(1404.58,472.97)(0.492,-2.349){21}{\rule{0.119pt}{1.933pt}}
\multiput(1403.17,476.99)(12.000,-50.987){2}{\rule{0.400pt}{0.967pt}}
\multiput(1416.58,420.60)(0.492,-1.534){19}{\rule{0.118pt}{1.300pt}}
\multiput(1415.17,423.30)(11.000,-30.302){2}{\rule{0.400pt}{0.650pt}}
\multiput(1427.58,389.82)(0.492,-0.841){21}{\rule{0.119pt}{0.767pt}}
\multiput(1426.17,391.41)(12.000,-18.409){2}{\rule{0.400pt}{0.383pt}}
\put(746.0,340.0){\rule[-0.200pt]{2.650pt}{0.400pt}}
\put(191.0,131.0){\rule[-0.200pt]{0.400pt}{155.380pt}}
\put(191.0,131.0){\rule[-0.200pt]{300.643pt}{0.400pt}}
\put(1439.0,131.0){\rule[-0.200pt]{0.400pt}{155.380pt}}
\put(191.0,776.0){\rule[-0.200pt]{300.643pt}{0.400pt}}
\end{picture}

\end{center}

\begin{center}
% GNUPLOT: LaTeX picture
\setlength{\unitlength}{0.240900pt}
\ifx\plotpoint\undefined\newsavebox{\plotpoint}\fi
\sbox{\plotpoint}{\rule[-0.200pt]{0.400pt}{0.400pt}}%
\begin{picture}(1500,900)(0,0)
\sbox{\plotpoint}{\rule[-0.200pt]{0.400pt}{0.400pt}}%
\put(171.0,131.0){\rule[-0.200pt]{4.818pt}{0.400pt}}
\put(151,131){\makebox(0,0)[r]{$0$}}
\put(1419.0,131.0){\rule[-0.200pt]{4.818pt}{0.400pt}}
\put(171.0,238.0){\rule[-0.200pt]{4.818pt}{0.400pt}}
\put(151,238){\makebox(0,0)[r]{$0.05$}}
\put(1419.0,238.0){\rule[-0.200pt]{4.818pt}{0.400pt}}
\put(171.0,346.0){\rule[-0.200pt]{4.818pt}{0.400pt}}
\put(151,346){\makebox(0,0)[r]{$0.1$}}
\put(1419.0,346.0){\rule[-0.200pt]{4.818pt}{0.400pt}}
\put(171.0,453.0){\rule[-0.200pt]{4.818pt}{0.400pt}}
\put(151,453){\makebox(0,0)[r]{$0.15$}}
\put(1419.0,453.0){\rule[-0.200pt]{4.818pt}{0.400pt}}
\put(171.0,561.0){\rule[-0.200pt]{4.818pt}{0.400pt}}
\put(151,561){\makebox(0,0)[r]{$0.2$}}
\put(1419.0,561.0){\rule[-0.200pt]{4.818pt}{0.400pt}}
\put(171.0,668.0){\rule[-0.200pt]{4.818pt}{0.400pt}}
\put(151,668){\makebox(0,0)[r]{$0.25$}}
\put(1419.0,668.0){\rule[-0.200pt]{4.818pt}{0.400pt}}
\put(171.0,776.0){\rule[-0.200pt]{4.818pt}{0.400pt}}
\put(151,776){\makebox(0,0)[r]{$0.3$}}
\put(1419.0,776.0){\rule[-0.200pt]{4.818pt}{0.400pt}}
\put(171.0,131.0){\rule[-0.200pt]{0.400pt}{4.818pt}}
\put(171,90){\makebox(0,0){$0$}}
\put(171.0,756.0){\rule[-0.200pt]{0.400pt}{4.818pt}}
\put(298.0,131.0){\rule[-0.200pt]{0.400pt}{4.818pt}}
\put(298,90){\makebox(0,0){$0.05$}}
\put(298.0,756.0){\rule[-0.200pt]{0.400pt}{4.818pt}}
\put(425.0,131.0){\rule[-0.200pt]{0.400pt}{4.818pt}}
\put(425,90){\makebox(0,0){$0.1$}}
\put(425.0,756.0){\rule[-0.200pt]{0.400pt}{4.818pt}}
\put(551.0,131.0){\rule[-0.200pt]{0.400pt}{4.818pt}}
\put(551,90){\makebox(0,0){$0.15$}}
\put(551.0,756.0){\rule[-0.200pt]{0.400pt}{4.818pt}}
\put(678.0,131.0){\rule[-0.200pt]{0.400pt}{4.818pt}}
\put(678,90){\makebox(0,0){$0.2$}}
\put(678.0,756.0){\rule[-0.200pt]{0.400pt}{4.818pt}}
\put(805.0,131.0){\rule[-0.200pt]{0.400pt}{4.818pt}}
\put(805,90){\makebox(0,0){$0.25$}}
\put(805.0,756.0){\rule[-0.200pt]{0.400pt}{4.818pt}}
\put(932.0,131.0){\rule[-0.200pt]{0.400pt}{4.818pt}}
\put(932,90){\makebox(0,0){$0.3$}}
\put(932.0,756.0){\rule[-0.200pt]{0.400pt}{4.818pt}}
\put(1059.0,131.0){\rule[-0.200pt]{0.400pt}{4.818pt}}
\put(1059,90){\makebox(0,0){$0.35$}}
\put(1059.0,756.0){\rule[-0.200pt]{0.400pt}{4.818pt}}
\put(1185.0,131.0){\rule[-0.200pt]{0.400pt}{4.818pt}}
\put(1185,90){\makebox(0,0){$0.4$}}
\put(1185.0,756.0){\rule[-0.200pt]{0.400pt}{4.818pt}}
\put(1312.0,131.0){\rule[-0.200pt]{0.400pt}{4.818pt}}
\put(1312,90){\makebox(0,0){$0.45$}}
\put(1312.0,756.0){\rule[-0.200pt]{0.400pt}{4.818pt}}
\put(1439.0,131.0){\rule[-0.200pt]{0.400pt}{4.818pt}}
\put(1439,90){\makebox(0,0){$0.5$}}
\put(1439.0,756.0){\rule[-0.200pt]{0.400pt}{4.818pt}}
\put(171.0,131.0){\rule[-0.200pt]{0.400pt}{155.380pt}}
\put(171.0,131.0){\rule[-0.200pt]{305.461pt}{0.400pt}}
\put(1439.0,131.0){\rule[-0.200pt]{0.400pt}{155.380pt}}
\put(171.0,776.0){\rule[-0.200pt]{305.461pt}{0.400pt}}
\put(30,453){\makebox(0,0){$I/I_0$}}
\put(805,29){\makebox(0,0){$Delta$}}
\put(805,838){\makebox(0,0){$I/I_0 (Delta)$}}
\put(171,131){\usebox{\plotpoint}}
\multiput(171.00,131.59)(1.378,0.477){7}{\rule{1.140pt}{0.115pt}}
\multiput(171.00,130.17)(10.634,5.000){2}{\rule{0.570pt}{0.400pt}}
\multiput(184.00,136.59)(1.123,0.482){9}{\rule{0.967pt}{0.116pt}}
\multiput(184.00,135.17)(10.994,6.000){2}{\rule{0.483pt}{0.400pt}}
\multiput(197.00,142.59)(1.267,0.477){7}{\rule{1.060pt}{0.115pt}}
\multiput(197.00,141.17)(9.800,5.000){2}{\rule{0.530pt}{0.400pt}}
\multiput(209.00,147.59)(1.123,0.482){9}{\rule{0.967pt}{0.116pt}}
\multiput(209.00,146.17)(10.994,6.000){2}{\rule{0.483pt}{0.400pt}}
\multiput(222.00,153.59)(1.378,0.477){7}{\rule{1.140pt}{0.115pt}}
\multiput(222.00,152.17)(10.634,5.000){2}{\rule{0.570pt}{0.400pt}}
\multiput(235.00,158.59)(1.123,0.482){9}{\rule{0.967pt}{0.116pt}}
\multiput(235.00,157.17)(10.994,6.000){2}{\rule{0.483pt}{0.400pt}}
\multiput(248.00,164.59)(1.378,0.477){7}{\rule{1.140pt}{0.115pt}}
\multiput(248.00,163.17)(10.634,5.000){2}{\rule{0.570pt}{0.400pt}}
\multiput(261.00,169.59)(1.267,0.477){7}{\rule{1.060pt}{0.115pt}}
\multiput(261.00,168.17)(9.800,5.000){2}{\rule{0.530pt}{0.400pt}}
\multiput(273.00,174.59)(1.123,0.482){9}{\rule{0.967pt}{0.116pt}}
\multiput(273.00,173.17)(10.994,6.000){2}{\rule{0.483pt}{0.400pt}}
\multiput(286.00,180.59)(1.378,0.477){7}{\rule{1.140pt}{0.115pt}}
\multiput(286.00,179.17)(10.634,5.000){2}{\rule{0.570pt}{0.400pt}}
\multiput(299.00,185.59)(1.123,0.482){9}{\rule{0.967pt}{0.116pt}}
\multiput(299.00,184.17)(10.994,6.000){2}{\rule{0.483pt}{0.400pt}}
\multiput(312.00,191.59)(1.378,0.477){7}{\rule{1.140pt}{0.115pt}}
\multiput(312.00,190.17)(10.634,5.000){2}{\rule{0.570pt}{0.400pt}}
\multiput(325.00,196.59)(1.123,0.482){9}{\rule{0.967pt}{0.116pt}}
\multiput(325.00,195.17)(10.994,6.000){2}{\rule{0.483pt}{0.400pt}}
\multiput(338.00,202.59)(1.267,0.477){7}{\rule{1.060pt}{0.115pt}}
\multiput(338.00,201.17)(9.800,5.000){2}{\rule{0.530pt}{0.400pt}}
\multiput(350.00,207.59)(1.378,0.477){7}{\rule{1.140pt}{0.115pt}}
\multiput(350.00,206.17)(10.634,5.000){2}{\rule{0.570pt}{0.400pt}}
\multiput(363.00,212.59)(1.123,0.482){9}{\rule{0.967pt}{0.116pt}}
\multiput(363.00,211.17)(10.994,6.000){2}{\rule{0.483pt}{0.400pt}}
\multiput(376.00,218.59)(1.378,0.477){7}{\rule{1.140pt}{0.115pt}}
\multiput(376.00,217.17)(10.634,5.000){2}{\rule{0.570pt}{0.400pt}}
\multiput(389.00,223.59)(1.123,0.482){9}{\rule{0.967pt}{0.116pt}}
\multiput(389.00,222.17)(10.994,6.000){2}{\rule{0.483pt}{0.400pt}}
\multiput(402.00,229.59)(1.267,0.477){7}{\rule{1.060pt}{0.115pt}}
\multiput(402.00,228.17)(9.800,5.000){2}{\rule{0.530pt}{0.400pt}}
\multiput(414.00,234.59)(1.123,0.482){9}{\rule{0.967pt}{0.116pt}}
\multiput(414.00,233.17)(10.994,6.000){2}{\rule{0.483pt}{0.400pt}}
\multiput(427.00,240.59)(1.378,0.477){7}{\rule{1.140pt}{0.115pt}}
\multiput(427.00,239.17)(10.634,5.000){2}{\rule{0.570pt}{0.400pt}}
\multiput(440.00,245.59)(1.378,0.477){7}{\rule{1.140pt}{0.115pt}}
\multiput(440.00,244.17)(10.634,5.000){2}{\rule{0.570pt}{0.400pt}}
\multiput(453.00,250.59)(1.123,0.482){9}{\rule{0.967pt}{0.116pt}}
\multiput(453.00,249.17)(10.994,6.000){2}{\rule{0.483pt}{0.400pt}}
\multiput(466.00,256.59)(1.267,0.477){7}{\rule{1.060pt}{0.115pt}}
\multiput(466.00,255.17)(9.800,5.000){2}{\rule{0.530pt}{0.400pt}}
\multiput(478.00,261.59)(1.123,0.482){9}{\rule{0.967pt}{0.116pt}}
\multiput(478.00,260.17)(10.994,6.000){2}{\rule{0.483pt}{0.400pt}}
\multiput(491.00,267.59)(1.378,0.477){7}{\rule{1.140pt}{0.115pt}}
\multiput(491.00,266.17)(10.634,5.000){2}{\rule{0.570pt}{0.400pt}}
\multiput(504.00,272.59)(1.123,0.482){9}{\rule{0.967pt}{0.116pt}}
\multiput(504.00,271.17)(10.994,6.000){2}{\rule{0.483pt}{0.400pt}}
\multiput(517.00,278.59)(1.378,0.477){7}{\rule{1.140pt}{0.115pt}}
\multiput(517.00,277.17)(10.634,5.000){2}{\rule{0.570pt}{0.400pt}}
\multiput(530.00,283.59)(1.267,0.477){7}{\rule{1.060pt}{0.115pt}}
\multiput(530.00,282.17)(9.800,5.000){2}{\rule{0.530pt}{0.400pt}}
\multiput(542.00,288.59)(1.123,0.482){9}{\rule{0.967pt}{0.116pt}}
\multiput(542.00,287.17)(10.994,6.000){2}{\rule{0.483pt}{0.400pt}}
\multiput(555.00,294.59)(1.378,0.477){7}{\rule{1.140pt}{0.115pt}}
\multiput(555.00,293.17)(10.634,5.000){2}{\rule{0.570pt}{0.400pt}}
\multiput(568.00,299.59)(1.123,0.482){9}{\rule{0.967pt}{0.116pt}}
\multiput(568.00,298.17)(10.994,6.000){2}{\rule{0.483pt}{0.400pt}}
\multiput(581.00,305.59)(1.378,0.477){7}{\rule{1.140pt}{0.115pt}}
\multiput(581.00,304.17)(10.634,5.000){2}{\rule{0.570pt}{0.400pt}}
\multiput(594.00,310.59)(1.033,0.482){9}{\rule{0.900pt}{0.116pt}}
\multiput(594.00,309.17)(10.132,6.000){2}{\rule{0.450pt}{0.400pt}}
\multiput(606.00,316.59)(1.378,0.477){7}{\rule{1.140pt}{0.115pt}}
\multiput(606.00,315.17)(10.634,5.000){2}{\rule{0.570pt}{0.400pt}}
\multiput(619.00,321.59)(1.378,0.477){7}{\rule{1.140pt}{0.115pt}}
\multiput(619.00,320.17)(10.634,5.000){2}{\rule{0.570pt}{0.400pt}}
\multiput(632.00,326.59)(1.123,0.482){9}{\rule{0.967pt}{0.116pt}}
\multiput(632.00,325.17)(10.994,6.000){2}{\rule{0.483pt}{0.400pt}}
\multiput(645.00,332.59)(1.378,0.477){7}{\rule{1.140pt}{0.115pt}}
\multiput(645.00,331.17)(10.634,5.000){2}{\rule{0.570pt}{0.400pt}}
\multiput(658.00,337.59)(1.123,0.482){9}{\rule{0.967pt}{0.116pt}}
\multiput(658.00,336.17)(10.994,6.000){2}{\rule{0.483pt}{0.400pt}}
\multiput(671.00,343.59)(1.267,0.477){7}{\rule{1.060pt}{0.115pt}}
\multiput(671.00,342.17)(9.800,5.000){2}{\rule{0.530pt}{0.400pt}}
\multiput(683.00,348.59)(1.123,0.482){9}{\rule{0.967pt}{0.116pt}}
\multiput(683.00,347.17)(10.994,6.000){2}{\rule{0.483pt}{0.400pt}}
\multiput(696.00,354.59)(1.378,0.477){7}{\rule{1.140pt}{0.115pt}}
\multiput(696.00,353.17)(10.634,5.000){2}{\rule{0.570pt}{0.400pt}}
\multiput(709.00,359.59)(1.378,0.477){7}{\rule{1.140pt}{0.115pt}}
\multiput(709.00,358.17)(10.634,5.000){2}{\rule{0.570pt}{0.400pt}}
\multiput(722.00,364.59)(1.123,0.482){9}{\rule{0.967pt}{0.116pt}}
\multiput(722.00,363.17)(10.994,6.000){2}{\rule{0.483pt}{0.400pt}}
\multiput(735.00,370.59)(1.267,0.477){7}{\rule{1.060pt}{0.115pt}}
\multiput(735.00,369.17)(9.800,5.000){2}{\rule{0.530pt}{0.400pt}}
\multiput(747.00,375.59)(1.123,0.482){9}{\rule{0.967pt}{0.116pt}}
\multiput(747.00,374.17)(10.994,6.000){2}{\rule{0.483pt}{0.400pt}}
\multiput(760.00,381.59)(1.378,0.477){7}{\rule{1.140pt}{0.115pt}}
\multiput(760.00,380.17)(10.634,5.000){2}{\rule{0.570pt}{0.400pt}}
\multiput(773.00,386.59)(1.123,0.482){9}{\rule{0.967pt}{0.116pt}}
\multiput(773.00,385.17)(10.994,6.000){2}{\rule{0.483pt}{0.400pt}}
\multiput(786.00,392.59)(1.378,0.477){7}{\rule{1.140pt}{0.115pt}}
\multiput(786.00,391.17)(10.634,5.000){2}{\rule{0.570pt}{0.400pt}}
\multiput(799.00,397.59)(1.267,0.477){7}{\rule{1.060pt}{0.115pt}}
\multiput(799.00,396.17)(9.800,5.000){2}{\rule{0.530pt}{0.400pt}}
\multiput(811.00,402.59)(1.123,0.482){9}{\rule{0.967pt}{0.116pt}}
\multiput(811.00,401.17)(10.994,6.000){2}{\rule{0.483pt}{0.400pt}}
\multiput(824.00,408.59)(1.378,0.477){7}{\rule{1.140pt}{0.115pt}}
\multiput(824.00,407.17)(10.634,5.000){2}{\rule{0.570pt}{0.400pt}}
\multiput(837.00,413.59)(1.123,0.482){9}{\rule{0.967pt}{0.116pt}}
\multiput(837.00,412.17)(10.994,6.000){2}{\rule{0.483pt}{0.400pt}}
\multiput(850.00,419.59)(1.378,0.477){7}{\rule{1.140pt}{0.115pt}}
\multiput(850.00,418.17)(10.634,5.000){2}{\rule{0.570pt}{0.400pt}}
\multiput(863.00,424.59)(1.267,0.477){7}{\rule{1.060pt}{0.115pt}}
\multiput(863.00,423.17)(9.800,5.000){2}{\rule{0.530pt}{0.400pt}}
\multiput(875.00,429.59)(1.123,0.482){9}{\rule{0.967pt}{0.116pt}}
\multiput(875.00,428.17)(10.994,6.000){2}{\rule{0.483pt}{0.400pt}}
\multiput(888.00,435.59)(1.378,0.477){7}{\rule{1.140pt}{0.115pt}}
\multiput(888.00,434.17)(10.634,5.000){2}{\rule{0.570pt}{0.400pt}}
\multiput(901.00,440.59)(1.123,0.482){9}{\rule{0.967pt}{0.116pt}}
\multiput(901.00,439.17)(10.994,6.000){2}{\rule{0.483pt}{0.400pt}}
\multiput(914.00,446.59)(1.378,0.477){7}{\rule{1.140pt}{0.115pt}}
\multiput(914.00,445.17)(10.634,5.000){2}{\rule{0.570pt}{0.400pt}}
\multiput(927.00,451.59)(1.033,0.482){9}{\rule{0.900pt}{0.116pt}}
\multiput(927.00,450.17)(10.132,6.000){2}{\rule{0.450pt}{0.400pt}}
\multiput(939.00,457.59)(1.378,0.477){7}{\rule{1.140pt}{0.115pt}}
\multiput(939.00,456.17)(10.634,5.000){2}{\rule{0.570pt}{0.400pt}}
\multiput(952.00,462.59)(1.123,0.482){9}{\rule{0.967pt}{0.116pt}}
\multiput(952.00,461.17)(10.994,6.000){2}{\rule{0.483pt}{0.400pt}}
\multiput(965.00,468.59)(1.378,0.477){7}{\rule{1.140pt}{0.115pt}}
\multiput(965.00,467.17)(10.634,5.000){2}{\rule{0.570pt}{0.400pt}}
\multiput(978.00,473.59)(1.123,0.482){9}{\rule{0.967pt}{0.116pt}}
\multiput(978.00,472.17)(10.994,6.000){2}{\rule{0.483pt}{0.400pt}}
\multiput(991.00,479.59)(1.378,0.477){7}{\rule{1.140pt}{0.115pt}}
\multiput(991.00,478.17)(10.634,5.000){2}{\rule{0.570pt}{0.400pt}}
\multiput(1004.00,484.59)(1.267,0.477){7}{\rule{1.060pt}{0.115pt}}
\multiput(1004.00,483.17)(9.800,5.000){2}{\rule{0.530pt}{0.400pt}}
\multiput(1016.00,489.59)(1.123,0.482){9}{\rule{0.967pt}{0.116pt}}
\multiput(1016.00,488.17)(10.994,6.000){2}{\rule{0.483pt}{0.400pt}}
\multiput(1029.00,495.59)(1.378,0.477){7}{\rule{1.140pt}{0.115pt}}
\multiput(1029.00,494.17)(10.634,5.000){2}{\rule{0.570pt}{0.400pt}}
\multiput(1042.00,500.59)(1.378,0.477){7}{\rule{1.140pt}{0.115pt}}
\multiput(1042.00,499.17)(10.634,5.000){2}{\rule{0.570pt}{0.400pt}}
\multiput(1055.00,505.59)(1.123,0.482){9}{\rule{0.967pt}{0.116pt}}
\multiput(1055.00,504.17)(10.994,6.000){2}{\rule{0.483pt}{0.400pt}}
\multiput(1068.00,511.59)(1.267,0.477){7}{\rule{1.060pt}{0.115pt}}
\multiput(1068.00,510.17)(9.800,5.000){2}{\rule{0.530pt}{0.400pt}}
\multiput(1080.00,516.59)(1.123,0.482){9}{\rule{0.967pt}{0.116pt}}
\multiput(1080.00,515.17)(10.994,6.000){2}{\rule{0.483pt}{0.400pt}}
\multiput(1093.00,522.59)(0.950,0.485){11}{\rule{0.843pt}{0.117pt}}
\multiput(1093.00,521.17)(11.251,7.000){2}{\rule{0.421pt}{0.400pt}}
\multiput(1106.00,529.59)(1.123,0.482){9}{\rule{0.967pt}{0.116pt}}
\multiput(1106.00,528.17)(10.994,6.000){2}{\rule{0.483pt}{0.400pt}}
\multiput(1119.00,535.59)(0.950,0.485){11}{\rule{0.843pt}{0.117pt}}
\multiput(1119.00,534.17)(11.251,7.000){2}{\rule{0.421pt}{0.400pt}}
\multiput(1132.00,542.59)(1.033,0.482){9}{\rule{0.900pt}{0.116pt}}
\multiput(1132.00,541.17)(10.132,6.000){2}{\rule{0.450pt}{0.400pt}}
\multiput(1144.00,548.59)(0.950,0.485){11}{\rule{0.843pt}{0.117pt}}
\multiput(1144.00,547.17)(11.251,7.000){2}{\rule{0.421pt}{0.400pt}}
\multiput(1157.00,555.59)(1.123,0.482){9}{\rule{0.967pt}{0.116pt}}
\multiput(1157.00,554.17)(10.994,6.000){2}{\rule{0.483pt}{0.400pt}}
\multiput(1170.00,561.59)(1.378,0.477){7}{\rule{1.140pt}{0.115pt}}
\multiput(1170.00,560.17)(10.634,5.000){2}{\rule{0.570pt}{0.400pt}}
\multiput(1183.00,566.59)(1.378,0.477){7}{\rule{1.140pt}{0.115pt}}
\multiput(1183.00,565.17)(10.634,5.000){2}{\rule{0.570pt}{0.400pt}}
\multiput(1196.00,571.59)(1.267,0.477){7}{\rule{1.060pt}{0.115pt}}
\multiput(1196.00,570.17)(9.800,5.000){2}{\rule{0.530pt}{0.400pt}}
\multiput(1208.00,576.61)(2.695,0.447){3}{\rule{1.833pt}{0.108pt}}
\multiput(1208.00,575.17)(9.195,3.000){2}{\rule{0.917pt}{0.400pt}}
\multiput(1221.00,579.61)(2.695,0.447){3}{\rule{1.833pt}{0.108pt}}
\multiput(1221.00,578.17)(9.195,3.000){2}{\rule{0.917pt}{0.400pt}}
\put(1234,582.17){\rule{2.700pt}{0.400pt}}
\multiput(1234.00,581.17)(7.396,2.000){2}{\rule{1.350pt}{0.400pt}}
\put(1247,584.17){\rule{2.700pt}{0.400pt}}
\multiput(1247.00,583.17)(7.396,2.000){2}{\rule{1.350pt}{0.400pt}}
\put(1260,586.17){\rule{2.500pt}{0.400pt}}
\multiput(1260.00,585.17)(6.811,2.000){2}{\rule{1.250pt}{0.400pt}}
\multiput(1272.00,588.59)(1.378,0.477){7}{\rule{1.140pt}{0.115pt}}
\multiput(1272.00,587.17)(10.634,5.000){2}{\rule{0.570pt}{0.400pt}}
\multiput(1285.00,593.59)(0.950,0.485){11}{\rule{0.843pt}{0.117pt}}
\multiput(1285.00,592.17)(11.251,7.000){2}{\rule{0.421pt}{0.400pt}}
\multiput(1298.00,600.59)(0.824,0.488){13}{\rule{0.750pt}{0.117pt}}
\multiput(1298.00,599.17)(11.443,8.000){2}{\rule{0.375pt}{0.400pt}}
\multiput(1311.00,608.59)(0.728,0.489){15}{\rule{0.678pt}{0.118pt}}
\multiput(1311.00,607.17)(11.593,9.000){2}{\rule{0.339pt}{0.400pt}}
\multiput(1324.00,617.59)(0.824,0.488){13}{\rule{0.750pt}{0.117pt}}
\multiput(1324.00,616.17)(11.443,8.000){2}{\rule{0.375pt}{0.400pt}}
\multiput(1337.00,625.59)(1.267,0.477){7}{\rule{1.060pt}{0.115pt}}
\multiput(1337.00,624.17)(9.800,5.000){2}{\rule{0.530pt}{0.400pt}}
\multiput(1349.00,630.61)(2.695,0.447){3}{\rule{1.833pt}{0.108pt}}
\multiput(1349.00,629.17)(9.195,3.000){2}{\rule{0.917pt}{0.400pt}}
\multiput(1362.00,633.61)(2.695,0.447){3}{\rule{1.833pt}{0.108pt}}
\multiput(1362.00,632.17)(9.195,3.000){2}{\rule{0.917pt}{0.400pt}}
\multiput(1375.00,636.60)(1.797,0.468){5}{\rule{1.400pt}{0.113pt}}
\multiput(1375.00,635.17)(10.094,4.000){2}{\rule{0.700pt}{0.400pt}}
\multiput(1388.00,640.59)(1.378,0.477){7}{\rule{1.140pt}{0.115pt}}
\multiput(1388.00,639.17)(10.634,5.000){2}{\rule{0.570pt}{0.400pt}}
\multiput(1401.00,645.59)(0.874,0.485){11}{\rule{0.786pt}{0.117pt}}
\multiput(1401.00,644.17)(10.369,7.000){2}{\rule{0.393pt}{0.400pt}}
\multiput(1413.00,652.59)(0.824,0.488){13}{\rule{0.750pt}{0.117pt}}
\multiput(1413.00,651.17)(11.443,8.000){2}{\rule{0.375pt}{0.400pt}}
\multiput(1426.00,660.59)(0.728,0.489){15}{\rule{0.678pt}{0.118pt}}
\multiput(1426.00,659.17)(11.593,9.000){2}{\rule{0.339pt}{0.400pt}}
\put(171.0,131.0){\rule[-0.200pt]{0.400pt}{155.380pt}}
\put(171.0,131.0){\rule[-0.200pt]{305.461pt}{0.400pt}}
\put(1439.0,131.0){\rule[-0.200pt]{0.400pt}{155.380pt}}
\put(171.0,776.0){\rule[-0.200pt]{305.461pt}{0.400pt}}
\end{picture}

\end{center}

\begin{center}
% GNUPLOT: LaTeX picture
\setlength{\unitlength}{0.240900pt}
\ifx\plotpoint\undefined\newsavebox{\plotpoint}\fi
\sbox{\plotpoint}{\rule[-0.200pt]{0.400pt}{0.400pt}}%
\begin{picture}(1500,900)(0,0)
\sbox{\plotpoint}{\rule[-0.200pt]{0.400pt}{0.400pt}}%
\put(231.0,131.0){\rule[-0.200pt]{4.818pt}{0.400pt}}
\put(211,131){\makebox(0,0)[r]{$0$}}
\put(1419.0,131.0){\rule[-0.200pt]{4.818pt}{0.400pt}}
\put(231.0,212.0){\rule[-0.200pt]{4.818pt}{0.400pt}}
\put(211,212){\makebox(0,0)[r]{$200000$}}
\put(1419.0,212.0){\rule[-0.200pt]{4.818pt}{0.400pt}}
\put(231.0,292.0){\rule[-0.200pt]{4.818pt}{0.400pt}}
\put(211,292){\makebox(0,0)[r]{$400000$}}
\put(1419.0,292.0){\rule[-0.200pt]{4.818pt}{0.400pt}}
\put(231.0,373.0){\rule[-0.200pt]{4.818pt}{0.400pt}}
\put(211,373){\makebox(0,0)[r]{$600000$}}
\put(1419.0,373.0){\rule[-0.200pt]{4.818pt}{0.400pt}}
\put(231.0,454.0){\rule[-0.200pt]{4.818pt}{0.400pt}}
\put(211,454){\makebox(0,0)[r]{$800000$}}
\put(1419.0,454.0){\rule[-0.200pt]{4.818pt}{0.400pt}}
\put(231.0,534.0){\rule[-0.200pt]{4.818pt}{0.400pt}}
\put(211,534){\makebox(0,0)[r]{$1e+06$}}
\put(1419.0,534.0){\rule[-0.200pt]{4.818pt}{0.400pt}}
\put(231.0,615.0){\rule[-0.200pt]{4.818pt}{0.400pt}}
\put(211,615){\makebox(0,0)[r]{$1.2e+06$}}
\put(1419.0,615.0){\rule[-0.200pt]{4.818pt}{0.400pt}}
\put(231.0,695.0){\rule[-0.200pt]{4.818pt}{0.400pt}}
\put(211,695){\makebox(0,0)[r]{$1.4e+06$}}
\put(1419.0,695.0){\rule[-0.200pt]{4.818pt}{0.400pt}}
\put(231.0,776.0){\rule[-0.200pt]{4.818pt}{0.400pt}}
\put(211,776){\makebox(0,0)[r]{$1.6e+06$}}
\put(1419.0,776.0){\rule[-0.200pt]{4.818pt}{0.400pt}}
\put(231.0,131.0){\rule[-0.200pt]{0.400pt}{4.818pt}}
\put(231,90){\makebox(0,0){$0$}}
\put(231.0,756.0){\rule[-0.200pt]{0.400pt}{4.818pt}}
\put(352.0,131.0){\rule[-0.200pt]{0.400pt}{4.818pt}}
\put(352,90){\makebox(0,0){$0.05$}}
\put(352.0,756.0){\rule[-0.200pt]{0.400pt}{4.818pt}}
\put(473.0,131.0){\rule[-0.200pt]{0.400pt}{4.818pt}}
\put(473,90){\makebox(0,0){$0.1$}}
\put(473.0,756.0){\rule[-0.200pt]{0.400pt}{4.818pt}}
\put(593.0,131.0){\rule[-0.200pt]{0.400pt}{4.818pt}}
\put(593,90){\makebox(0,0){$0.15$}}
\put(593.0,756.0){\rule[-0.200pt]{0.400pt}{4.818pt}}
\put(714.0,131.0){\rule[-0.200pt]{0.400pt}{4.818pt}}
\put(714,90){\makebox(0,0){$0.2$}}
\put(714.0,756.0){\rule[-0.200pt]{0.400pt}{4.818pt}}
\put(835.0,131.0){\rule[-0.200pt]{0.400pt}{4.818pt}}
\put(835,90){\makebox(0,0){$0.25$}}
\put(835.0,756.0){\rule[-0.200pt]{0.400pt}{4.818pt}}
\put(956.0,131.0){\rule[-0.200pt]{0.400pt}{4.818pt}}
\put(956,90){\makebox(0,0){$0.3$}}
\put(956.0,756.0){\rule[-0.200pt]{0.400pt}{4.818pt}}
\put(1077.0,131.0){\rule[-0.200pt]{0.400pt}{4.818pt}}
\put(1077,90){\makebox(0,0){$0.35$}}
\put(1077.0,756.0){\rule[-0.200pt]{0.400pt}{4.818pt}}
\put(1197.0,131.0){\rule[-0.200pt]{0.400pt}{4.818pt}}
\put(1197,90){\makebox(0,0){$0.4$}}
\put(1197.0,756.0){\rule[-0.200pt]{0.400pt}{4.818pt}}
\put(1318.0,131.0){\rule[-0.200pt]{0.400pt}{4.818pt}}
\put(1318,90){\makebox(0,0){$0.45$}}
\put(1318.0,756.0){\rule[-0.200pt]{0.400pt}{4.818pt}}
\put(1439.0,131.0){\rule[-0.200pt]{0.400pt}{4.818pt}}
\put(1439,90){\makebox(0,0){$0.5$}}
\put(1439.0,756.0){\rule[-0.200pt]{0.400pt}{4.818pt}}
\put(231.0,131.0){\rule[-0.200pt]{0.400pt}{155.380pt}}
\put(231.0,131.0){\rule[-0.200pt]{291.007pt}{0.400pt}}
\put(1439.0,131.0){\rule[-0.200pt]{0.400pt}{155.380pt}}
\put(231.0,776.0){\rule[-0.200pt]{291.007pt}{0.400pt}}
\put(30,453){\makebox(0,0){$U**2$}}
\put(835,29){\makebox(0,0){$D$}}
\put(835,838){\makebox(0,0){$D(U**2)$}}
\put(231,156){\usebox{\plotpoint}}
\put(231,154.17){\rule{2.500pt}{0.400pt}}
\multiput(231.00,155.17)(6.811,-2.000){2}{\rule{1.250pt}{0.400pt}}
\multiput(255.00,154.61)(2.695,0.447){3}{\rule{1.833pt}{0.108pt}}
\multiput(255.00,153.17)(9.195,3.000){2}{\rule{0.917pt}{0.400pt}}
\multiput(268.00,157.59)(0.758,0.488){13}{\rule{0.700pt}{0.117pt}}
\multiput(268.00,156.17)(10.547,8.000){2}{\rule{0.350pt}{0.400pt}}
\multiput(280.58,165.00)(0.492,0.582){21}{\rule{0.119pt}{0.567pt}}
\multiput(279.17,165.00)(12.000,12.824){2}{\rule{0.400pt}{0.283pt}}
\multiput(292.58,179.00)(0.492,0.669){21}{\rule{0.119pt}{0.633pt}}
\multiput(291.17,179.00)(12.000,14.685){2}{\rule{0.400pt}{0.317pt}}
\multiput(304.00,195.58)(0.496,0.492){21}{\rule{0.500pt}{0.119pt}}
\multiput(304.00,194.17)(10.962,12.000){2}{\rule{0.250pt}{0.400pt}}
\multiput(316.00,207.59)(1.123,0.482){9}{\rule{0.967pt}{0.116pt}}
\multiput(316.00,206.17)(10.994,6.000){2}{\rule{0.483pt}{0.400pt}}
\multiput(329.00,211.95)(2.472,-0.447){3}{\rule{1.700pt}{0.108pt}}
\multiput(329.00,212.17)(8.472,-3.000){2}{\rule{0.850pt}{0.400pt}}
\put(341,208.17){\rule{2.500pt}{0.400pt}}
\multiput(341.00,209.17)(6.811,-2.000){2}{\rule{1.250pt}{0.400pt}}
\multiput(353.58,208.00)(0.492,0.539){21}{\rule{0.119pt}{0.533pt}}
\multiput(352.17,208.00)(12.000,11.893){2}{\rule{0.400pt}{0.267pt}}
\multiput(365.58,221.00)(0.492,1.659){21}{\rule{0.119pt}{1.400pt}}
\multiput(364.17,221.00)(12.000,36.094){2}{\rule{0.400pt}{0.700pt}}
\multiput(377.58,260.00)(0.493,2.836){23}{\rule{0.119pt}{2.315pt}}
\multiput(376.17,260.00)(13.000,67.194){2}{\rule{0.400pt}{1.158pt}}
\multiput(390.58,332.00)(0.492,4.029){21}{\rule{0.119pt}{3.233pt}}
\multiput(389.17,332.00)(12.000,87.289){2}{\rule{0.400pt}{1.617pt}}
\multiput(402.58,426.00)(0.492,4.417){21}{\rule{0.119pt}{3.533pt}}
\multiput(401.17,426.00)(12.000,95.666){2}{\rule{0.400pt}{1.767pt}}
\multiput(414.58,529.00)(0.492,4.158){21}{\rule{0.119pt}{3.333pt}}
\multiput(413.17,529.00)(12.000,90.081){2}{\rule{0.400pt}{1.667pt}}
\multiput(426.58,626.00)(0.492,3.383){21}{\rule{0.119pt}{2.733pt}}
\multiput(425.17,626.00)(12.000,73.327){2}{\rule{0.400pt}{1.367pt}}
\multiput(438.58,705.00)(0.493,1.845){23}{\rule{0.119pt}{1.546pt}}
\multiput(437.17,705.00)(13.000,43.791){2}{\rule{0.400pt}{0.773pt}}
\put(451,751.67){\rule{2.891pt}{0.400pt}}
\multiput(451.00,751.17)(6.000,1.000){2}{\rule{1.445pt}{0.400pt}}
\multiput(463.58,744.42)(0.492,-2.521){21}{\rule{0.119pt}{2.067pt}}
\multiput(462.17,748.71)(12.000,-54.711){2}{\rule{0.400pt}{1.033pt}}
\multiput(475.58,677.67)(0.492,-4.934){21}{\rule{0.119pt}{3.933pt}}
\multiput(474.17,685.84)(12.000,-106.836){2}{\rule{0.400pt}{1.967pt}}
\multiput(487.58,559.90)(0.492,-5.795){21}{\rule{0.119pt}{4.600pt}}
\multiput(486.17,569.45)(12.000,-125.452){2}{\rule{0.400pt}{2.300pt}}
\multiput(499.58,429.28)(0.493,-4.422){23}{\rule{0.119pt}{3.546pt}}
\multiput(498.17,436.64)(13.000,-104.640){2}{\rule{0.400pt}{1.773pt}}
\multiput(512.58,325.08)(0.492,-2.004){21}{\rule{0.119pt}{1.667pt}}
\multiput(511.17,328.54)(12.000,-43.541){2}{\rule{0.400pt}{0.833pt}}
\multiput(524.58,285.00)(0.492,1.487){21}{\rule{0.119pt}{1.267pt}}
\multiput(523.17,285.00)(12.000,32.371){2}{\rule{0.400pt}{0.633pt}}
\multiput(536.58,320.00)(0.492,3.512){21}{\rule{0.119pt}{2.833pt}}
\multiput(535.17,320.00)(12.000,76.119){2}{\rule{0.400pt}{1.417pt}}
\multiput(548.58,402.00)(0.492,3.684){21}{\rule{0.119pt}{2.967pt}}
\multiput(547.17,402.00)(12.000,79.843){2}{\rule{0.400pt}{1.483pt}}
\multiput(560.58,488.00)(0.493,1.924){23}{\rule{0.119pt}{1.608pt}}
\multiput(559.17,488.00)(13.000,45.663){2}{\rule{0.400pt}{0.804pt}}
\multiput(573.00,535.93)(1.033,-0.482){9}{\rule{0.900pt}{0.116pt}}
\multiput(573.00,536.17)(10.132,-6.000){2}{\rule{0.450pt}{0.400pt}}
\multiput(585.58,525.19)(0.492,-1.659){21}{\rule{0.119pt}{1.400pt}}
\multiput(584.17,528.09)(12.000,-36.094){2}{\rule{0.400pt}{0.700pt}}
\multiput(597.58,485.08)(0.492,-2.004){21}{\rule{0.119pt}{1.667pt}}
\multiput(596.17,488.54)(12.000,-43.541){2}{\rule{0.400pt}{0.833pt}}
\multiput(609.58,440.71)(0.492,-1.186){21}{\rule{0.119pt}{1.033pt}}
\multiput(608.17,442.86)(12.000,-25.855){2}{\rule{0.400pt}{0.517pt}}
\put(621,417.17){\rule{2.700pt}{0.400pt}}
\multiput(621.00,416.17)(7.396,2.000){2}{\rule{1.350pt}{0.400pt}}
\multiput(634.58,419.00)(0.492,0.970){21}{\rule{0.119pt}{0.867pt}}
\multiput(633.17,419.00)(12.000,21.201){2}{\rule{0.400pt}{0.433pt}}
\multiput(646.58,442.00)(0.492,1.358){21}{\rule{0.119pt}{1.167pt}}
\multiput(645.17,442.00)(12.000,29.579){2}{\rule{0.400pt}{0.583pt}}
\multiput(658.58,474.00)(0.492,1.272){21}{\rule{0.119pt}{1.100pt}}
\multiput(657.17,474.00)(12.000,27.717){2}{\rule{0.400pt}{0.550pt}}
\multiput(670.58,504.00)(0.492,0.884){21}{\rule{0.119pt}{0.800pt}}
\multiput(669.17,504.00)(12.000,19.340){2}{\rule{0.400pt}{0.400pt}}
\multiput(682.00,525.58)(0.497,0.493){23}{\rule{0.500pt}{0.119pt}}
\multiput(682.00,524.17)(11.962,13.000){2}{\rule{0.250pt}{0.400pt}}
\multiput(695.00,538.59)(1.033,0.482){9}{\rule{0.900pt}{0.116pt}}
\multiput(695.00,537.17)(10.132,6.000){2}{\rule{0.450pt}{0.400pt}}
\put(707,543.67){\rule{2.891pt}{0.400pt}}
\multiput(707.00,543.17)(6.000,1.000){2}{\rule{1.445pt}{0.400pt}}
\multiput(719.00,543.94)(1.651,-0.468){5}{\rule{1.300pt}{0.113pt}}
\multiput(719.00,544.17)(9.302,-4.000){2}{\rule{0.650pt}{0.400pt}}
\multiput(731.00,539.93)(0.874,-0.485){11}{\rule{0.786pt}{0.117pt}}
\multiput(731.00,540.17)(10.369,-7.000){2}{\rule{0.393pt}{0.400pt}}
\multiput(743.00,532.93)(0.824,-0.488){13}{\rule{0.750pt}{0.117pt}}
\multiput(743.00,533.17)(11.443,-8.000){2}{\rule{0.375pt}{0.400pt}}
\multiput(756.00,524.93)(0.669,-0.489){15}{\rule{0.633pt}{0.118pt}}
\multiput(756.00,525.17)(10.685,-9.000){2}{\rule{0.317pt}{0.400pt}}
\multiput(768.00,515.93)(0.758,-0.488){13}{\rule{0.700pt}{0.117pt}}
\multiput(768.00,516.17)(10.547,-8.000){2}{\rule{0.350pt}{0.400pt}}
\multiput(780.00,507.93)(1.267,-0.477){7}{\rule{1.060pt}{0.115pt}}
\multiput(780.00,508.17)(9.800,-5.000){2}{\rule{0.530pt}{0.400pt}}
\put(792,502.17){\rule{2.500pt}{0.400pt}}
\multiput(792.00,503.17)(6.811,-2.000){2}{\rule{1.250pt}{0.400pt}}
\multiput(804.00,502.61)(2.695,0.447){3}{\rule{1.833pt}{0.108pt}}
\multiput(804.00,501.17)(9.195,3.000){2}{\rule{0.917pt}{0.400pt}}
\multiput(817.00,505.59)(0.874,0.485){11}{\rule{0.786pt}{0.117pt}}
\multiput(817.00,504.17)(10.369,7.000){2}{\rule{0.393pt}{0.400pt}}
\multiput(829.58,512.00)(0.492,0.539){21}{\rule{0.119pt}{0.533pt}}
\multiput(828.17,512.00)(12.000,11.893){2}{\rule{0.400pt}{0.267pt}}
\multiput(841.58,525.00)(0.492,0.625){21}{\rule{0.119pt}{0.600pt}}
\multiput(840.17,525.00)(12.000,13.755){2}{\rule{0.400pt}{0.300pt}}
\multiput(853.58,540.00)(0.493,0.695){23}{\rule{0.119pt}{0.654pt}}
\multiput(852.17,540.00)(13.000,16.643){2}{\rule{0.400pt}{0.327pt}}
\multiput(866.58,558.00)(0.492,0.841){21}{\rule{0.119pt}{0.767pt}}
\multiput(865.17,558.00)(12.000,18.409){2}{\rule{0.400pt}{0.383pt}}
\multiput(878.58,578.00)(0.492,0.841){21}{\rule{0.119pt}{0.767pt}}
\multiput(877.17,578.00)(12.000,18.409){2}{\rule{0.400pt}{0.383pt}}
\multiput(890.58,598.00)(0.492,0.841){21}{\rule{0.119pt}{0.767pt}}
\multiput(889.17,598.00)(12.000,18.409){2}{\rule{0.400pt}{0.383pt}}
\multiput(902.58,618.00)(0.492,0.798){21}{\rule{0.119pt}{0.733pt}}
\multiput(901.17,618.00)(12.000,17.478){2}{\rule{0.400pt}{0.367pt}}
\multiput(914.58,637.00)(0.493,0.695){23}{\rule{0.119pt}{0.654pt}}
\multiput(913.17,637.00)(13.000,16.643){2}{\rule{0.400pt}{0.327pt}}
\multiput(927.58,655.00)(0.492,0.582){21}{\rule{0.119pt}{0.567pt}}
\multiput(926.17,655.00)(12.000,12.824){2}{\rule{0.400pt}{0.283pt}}
\multiput(939.00,669.58)(0.543,0.492){19}{\rule{0.536pt}{0.118pt}}
\multiput(939.00,668.17)(10.887,11.000){2}{\rule{0.268pt}{0.400pt}}
\multiput(951.00,680.59)(0.874,0.485){11}{\rule{0.786pt}{0.117pt}}
\multiput(951.00,679.17)(10.369,7.000){2}{\rule{0.393pt}{0.400pt}}
\put(963,686.67){\rule{2.891pt}{0.400pt}}
\multiput(963.00,686.17)(6.000,1.000){2}{\rule{1.445pt}{0.400pt}}
\multiput(975.00,686.93)(1.378,-0.477){7}{\rule{1.140pt}{0.115pt}}
\multiput(975.00,687.17)(10.634,-5.000){2}{\rule{0.570pt}{0.400pt}}
\multiput(988.58,680.79)(0.492,-0.539){21}{\rule{0.119pt}{0.533pt}}
\multiput(987.17,681.89)(12.000,-11.893){2}{\rule{0.400pt}{0.267pt}}
\multiput(1000.58,666.82)(0.492,-0.841){21}{\rule{0.119pt}{0.767pt}}
\multiput(999.17,668.41)(12.000,-18.409){2}{\rule{0.400pt}{0.383pt}}
\multiput(1012.58,645.71)(0.492,-1.186){21}{\rule{0.119pt}{1.033pt}}
\multiput(1011.17,647.86)(12.000,-25.855){2}{\rule{0.400pt}{0.517pt}}
\multiput(1024.58,616.60)(0.492,-1.530){21}{\rule{0.119pt}{1.300pt}}
\multiput(1023.17,619.30)(12.000,-33.302){2}{\rule{0.400pt}{0.650pt}}
\multiput(1036.58,579.84)(0.493,-1.765){23}{\rule{0.119pt}{1.485pt}}
\multiput(1035.17,582.92)(13.000,-41.919){2}{\rule{0.400pt}{0.742pt}}
\multiput(1049.58,533.53)(0.492,-2.176){21}{\rule{0.119pt}{1.800pt}}
\multiput(1048.17,537.26)(12.000,-47.264){2}{\rule{0.400pt}{0.900pt}}
\multiput(1061.58,481.97)(0.492,-2.349){21}{\rule{0.119pt}{1.933pt}}
\multiput(1060.17,485.99)(12.000,-50.987){2}{\rule{0.400pt}{0.967pt}}
\multiput(1073.58,427.39)(0.492,-2.219){21}{\rule{0.119pt}{1.833pt}}
\multiput(1072.17,431.19)(12.000,-48.195){2}{\rule{0.400pt}{0.917pt}}
\multiput(1085.58,376.50)(0.492,-1.875){21}{\rule{0.119pt}{1.567pt}}
\multiput(1084.17,379.75)(12.000,-40.748){2}{\rule{0.400pt}{0.783pt}}
\multiput(1097.58,334.63)(0.493,-1.210){23}{\rule{0.119pt}{1.054pt}}
\multiput(1096.17,336.81)(13.000,-28.813){2}{\rule{0.400pt}{0.527pt}}
\multiput(1110.58,305.37)(0.492,-0.669){21}{\rule{0.119pt}{0.633pt}}
\multiput(1109.17,306.69)(12.000,-14.685){2}{\rule{0.400pt}{0.317pt}}
\multiput(1122.00,290.95)(2.472,-0.447){3}{\rule{1.700pt}{0.108pt}}
\multiput(1122.00,291.17)(8.472,-3.000){2}{\rule{0.850pt}{0.400pt}}
\multiput(1134.00,289.59)(1.033,0.482){9}{\rule{0.900pt}{0.116pt}}
\multiput(1134.00,288.17)(10.132,6.000){2}{\rule{0.450pt}{0.400pt}}
\multiput(1146.58,295.00)(0.492,0.582){21}{\rule{0.119pt}{0.567pt}}
\multiput(1145.17,295.00)(12.000,12.824){2}{\rule{0.400pt}{0.283pt}}
\multiput(1158.58,309.00)(0.493,0.774){23}{\rule{0.119pt}{0.715pt}}
\multiput(1157.17,309.00)(13.000,18.515){2}{\rule{0.400pt}{0.358pt}}
\multiput(1171.58,329.00)(0.492,1.013){21}{\rule{0.119pt}{0.900pt}}
\multiput(1170.17,329.00)(12.000,22.132){2}{\rule{0.400pt}{0.450pt}}
\multiput(1183.58,353.00)(0.492,1.056){21}{\rule{0.119pt}{0.933pt}}
\multiput(1182.17,353.00)(12.000,23.063){2}{\rule{0.400pt}{0.467pt}}
\multiput(1195.58,378.00)(0.492,1.013){21}{\rule{0.119pt}{0.900pt}}
\multiput(1194.17,378.00)(12.000,22.132){2}{\rule{0.400pt}{0.450pt}}
\multiput(1207.58,402.00)(0.492,0.927){21}{\rule{0.119pt}{0.833pt}}
\multiput(1206.17,402.00)(12.000,20.270){2}{\rule{0.400pt}{0.417pt}}
\multiput(1219.58,424.00)(0.493,0.616){23}{\rule{0.119pt}{0.592pt}}
\multiput(1218.17,424.00)(13.000,14.771){2}{\rule{0.400pt}{0.296pt}}
\multiput(1232.00,440.59)(0.669,0.489){15}{\rule{0.633pt}{0.118pt}}
\multiput(1232.00,439.17)(10.685,9.000){2}{\rule{0.317pt}{0.400pt}}
\put(1244,448.67){\rule{2.891pt}{0.400pt}}
\multiput(1244.00,448.17)(6.000,1.000){2}{\rule{1.445pt}{0.400pt}}
\multiput(1256.00,448.94)(1.651,-0.468){5}{\rule{1.300pt}{0.113pt}}
\multiput(1256.00,449.17)(9.302,-4.000){2}{\rule{0.650pt}{0.400pt}}
\multiput(1268.00,444.95)(2.472,-0.447){3}{\rule{1.700pt}{0.108pt}}
\multiput(1268.00,445.17)(8.472,-3.000){2}{\rule{0.850pt}{0.400pt}}
\multiput(1280.00,443.60)(1.797,0.468){5}{\rule{1.400pt}{0.113pt}}
\multiput(1280.00,442.17)(10.094,4.000){2}{\rule{0.700pt}{0.400pt}}
\multiput(1293.58,447.00)(0.492,0.669){21}{\rule{0.119pt}{0.633pt}}
\multiput(1292.17,447.00)(12.000,14.685){2}{\rule{0.400pt}{0.317pt}}
\multiput(1305.58,463.00)(0.492,0.927){21}{\rule{0.119pt}{0.833pt}}
\multiput(1304.17,463.00)(12.000,20.270){2}{\rule{0.400pt}{0.417pt}}
\multiput(1317.58,485.00)(0.492,0.841){21}{\rule{0.119pt}{0.767pt}}
\multiput(1316.17,485.00)(12.000,18.409){2}{\rule{0.400pt}{0.383pt}}
\multiput(1329.00,505.59)(0.669,0.489){15}{\rule{0.633pt}{0.118pt}}
\multiput(1329.00,504.17)(10.685,9.000){2}{\rule{0.317pt}{0.400pt}}
\multiput(1341.00,512.93)(0.824,-0.488){13}{\rule{0.750pt}{0.117pt}}
\multiput(1341.00,513.17)(11.443,-8.000){2}{\rule{0.375pt}{0.400pt}}
\multiput(1354.58,502.54)(0.492,-0.927){21}{\rule{0.119pt}{0.833pt}}
\multiput(1353.17,504.27)(12.000,-20.270){2}{\rule{0.400pt}{0.417pt}}
\multiput(1366.58,479.71)(0.492,-1.186){21}{\rule{0.119pt}{1.033pt}}
\multiput(1365.17,481.86)(12.000,-25.855){2}{\rule{0.400pt}{0.517pt}}
\multiput(1378.58,451.71)(0.492,-1.186){21}{\rule{0.119pt}{1.033pt}}
\multiput(1377.17,453.86)(12.000,-25.855){2}{\rule{0.400pt}{0.517pt}}
\multiput(1390.58,424.68)(0.492,-0.884){21}{\rule{0.119pt}{0.800pt}}
\multiput(1389.17,426.34)(12.000,-19.340){2}{\rule{0.400pt}{0.400pt}}
\multiput(1402.58,404.67)(0.493,-0.576){23}{\rule{0.119pt}{0.562pt}}
\multiput(1401.17,405.83)(13.000,-13.834){2}{\rule{0.400pt}{0.281pt}}
\multiput(1415.00,390.92)(0.600,-0.491){17}{\rule{0.580pt}{0.118pt}}
\multiput(1415.00,391.17)(10.796,-10.000){2}{\rule{0.290pt}{0.400pt}}
\multiput(1427.00,380.93)(0.669,-0.489){15}{\rule{0.633pt}{0.118pt}}
\multiput(1427.00,381.17)(10.685,-9.000){2}{\rule{0.317pt}{0.400pt}}
\put(243.0,154.0){\rule[-0.200pt]{2.891pt}{0.400pt}}
\put(231.0,131.0){\rule[-0.200pt]{0.400pt}{155.380pt}}
\put(231.0,131.0){\rule[-0.200pt]{291.007pt}{0.400pt}}
\put(1439.0,131.0){\rule[-0.200pt]{0.400pt}{155.380pt}}
\put(231.0,776.0){\rule[-0.200pt]{291.007pt}{0.400pt}}
\end{picture}

\end{center}

\begin{table}[h!]
\begin{center}
\begin{tabular}{|c|c|c|c|}
\hline 
$U (Volts)$ & $I$ & $I/I_0$ & $\Delta$\\ 
\hline
100 & 7 & 0.01 & 0.01\\ 
\hline
125&8&0.01&0.01\\
\hline
150&6&0.00&0.00\\
\hline
175&5&0.00&0.00\\
\hline
200&3&0.00&0.00\\
\hline
225&9&0.01&0.01\\
\hline
250&5&0.00&0.00\\
\hline
275&1&0.00&0.00\\
\hline
300&0&0.00&0.00\\
\hline
325&2&0.00&0.00\\
\hline
350&7&0.01&0.01\\
\hline
375&11&0.01&0.01\\
\hline
400&16&0.01&0.01\\
\hline
425&21&0.02&0.02\\
\hline
450&26&0.02&0.02\\
\hline
475&30&0.02&0.02\\
\hline
500&37&0.03&0.03\\
\hline
525&38&0.03&0.03\\
\hline
550&38&0.03&0.03\\
\hline
575&32&0.03&0.03\\
\hline
600&40&0.03&0.03\\
\hline
625&77&0.06&0.06\\
\hline
650&112&0.09&0.09\\
\hline
675&157&0.12&0.12\\
\hline
700&232&0.18&0.18\\
\hline
725&270&0.21&0.22\\
\hline
750&297&0.23&0.24\\
\hline
775&310&0.25&0.25\\
\hline
800&303&0.24&0.24\\
\hline
825&265&0.21&0.21\\
\hline
850&97&0.08&0.08\\
\hline
875&90&0.07&0.07\\
\hline
900&150&0.12&0.12\\
\hline
925&212&0.17&0.17\\
\hline
950&263&0.21&0.21\\
\hline
975&290&0.23&0.23\\
\hline
1000&301&0.24&0.24\\
\hline
1025&280&0.22&0.22\\
\hline
1050&210&0.17&0.17\\
\hline
1075&152&0.12&0.12\\
\hline
1100&90&0.07&0.07\\
\hline
1125&148&0.12&0.12\\
\hline
1150&206&0.16&0.16\\
\hline
1175&113&0.09&0.09\\
\hline
1200&60&0.05&0.05\\
\hline
\end{tabular}
\end{center}
\caption{Τάση $V_{1}$, $V_{2}$ με $E_{B}=10V$ }
\end{table}

\section{Ερώτημα 2}


\subsection*{\underline{Λύση:}}


\section{Ερώτημα 3}


\subsection*{\underline{Λύση:}}


\end{document}
